% Intended LaTeX compiler: pdflatex
\documentclass[a4paper,12pt]{article}
\usepackage[utf8]{inputenc}
\usepackage[T1]{fontenc}
\usepackage{graphicx}
\usepackage{longtable}
\usepackage{wrapfig}
\usepackage{rotating}
\usepackage[normalem]{ulem}
\usepackage{amsmath}
\usepackage{amssymb}
\usepackage{capt-of}
\usepackage{hyperref}
\usepackage[margin=0.2in]{geometry}
\usepackage{amsmath}
\usepackage{amsfonts}
\usepackage{indentfirst}
\usepackage{pgffor}
\usepackage{amssymb}
\usepackage{cancel}
\usepackage{amsthm}
\usepackage{polynom}
\usepackage{tikz}
\usepackage[tikz]{bclogo}
\usepackage{listings}
\usepackage[most]{tcolorbox}
\usepackage{forest}
\usepackage{adjustbox}
\usepackage{tikz-3dplot}
\usepackage{mathtools}
\usepackage{pgfplots}
\usetikzlibrary{tikzmark,calc,fit,matrix,arrows,automata,positioning}
\usepackage{centernot}
\usepackage{lmodern}
\usepackage{physics}
\usepackage[boxed,linesnumbered,vlined]{algorithm2e}
\usepackage{algpseudocode}
\usepackage{algorithmicx}
\usepackage{svg}
\tikzstyle{every picture}+=[remember picture]
\DeclareMathOperator{\spn}{span}
\DeclareMathOperator{\dom}{domain}
\DeclareMathOperator{\ran}{range}
\DeclareMathOperator{\rng}{range}
\DeclareMathOperator{\img}{Im}
\DeclareMathOperator{\adj}{adj}
\newcommand\dif[1]{\:\textrm{d}#1}
\DeclarePairedDelimiter\ceil{\lceil}{\rceil}
\DeclarePairedDelimiter\floor{\lfloor}{\rfloor}
\newcommand{\ihat}{\hat{\textbf{\i}}}
\newcommand{\jhat}{\hat{\textbf{\j}}}
\newcommand{\khat}{\hat{\textbf{k}}}
\newcommand{\rhat}{\hat{\textbf{r}}}
\newcommand{\thetahat}{\boldsymbol{\hat{\theta}}}
\author{mahmood}
\date{\today}
\title{digital systems homework 3}
\hypersetup{
 pdfauthor={mahmood},
 pdftitle={digital systems homework 3},
 pdfkeywords={},
 pdfsubject={homework on the subject of },
 pdfcreator={Emacs 30.0.50 (Org mode 9.6)}, 
 pdflang={English}}
\begin{document}

\maketitle
\definecolor{Brown}{cmyk}{0,0.81,1,0.60}
\definecolor{OliveGreen}{cmyk}{0.64,0,0.95,0.40}
\definecolor{CadetBlue}{cmyk}{0.62,0.57,0.23,0}
\definecolor{lightlightgray}{gray}{0.9}
\lstset{
  language=C,                             % Code langugage
  basicstyle=\ttfamily,                   % Code font, Examples: \footnotesize, \ttfamily
  keywordstyle=\color{OliveGreen},        % Keywords font ('*' = uppercase)
  commentstyle=\color{gray},              % Comments font
  numbers=left,                           % Line nums position
  numberstyle=\tiny,                      % Line-numbers fonts
  stepnumber=1,                           % Step between two line-numbers
  numbersep=5pt,                          % How far are line-numbers from code
  backgroundcolor=\color{lightlightgray}, % Choose background color
  frame=single,                           % A frame around the code
  tabsize=2,                              % Default tab size
  captionpos=b,                           % Caption-position = bottom
  breaklines=true,                        % Automatic line breaking?
  breakatwhitespace=false,                % Automatic breaks only at whitespace?
  showspaces=false,                       % Dont make spaces visible
  showtabs=false,                         % Dont make tabls visible
  columns=flexible,                       % Column format
  morekeywords={__global__, __device__},  % CUDA specific keywords
  showstringspaces=false,                 % dont show spaces in strings
}
% so that latex doesnt complain about sage not being a language
\lstdefinelanguage{sage}{}

% environment definition for special blocks
\theoremstyle{definition}
\newtheorem{my_example}{Example}
\counterwithin{my_example}{section}
\counterwithin{my_example}{subsection}
\newtheorem{definition}{Definition}
\counterwithin{definition}{section}
\counterwithin{definition}{subsection}
\newtheorem{characteristic}{Characteristic}
\newtheorem{entailment}{Entailment}
%\newtheore*{proof}{Proof}
\newtheorem{lemma}{Lemma}
\newtheorem{question}{Question}
\newtheorem{subquestion}{Subquestion}[question]
%\newtheorem{note}{Note}
\newtheorem{answer}{Answer}
\counterwithin{answer}{question}
\counterwithin{answer}{subquestion}
\newtheorem{step}{Step}[answer]

% the following snippet auto resizes images to fit properly
% Determine if the image is too wide for the page.
% \makeatletter
% \def\ScaleIfNeeded{%
%   \ifdim\Gin@nat@width>\linewidth
%     \linewidth
%   \else
%     \Gin@nat@width
%   \fi
% }
% \makeatother
% % Resize figures that are too wide for the page.
% \let\oldincludegraphics\includegraphics
% \renewcommand\includegraphics[2][]{%
%   \oldincludegraphics[width=\ScaleIfNeeded]{#2}
% }

\newenvironment{note}
  {\begin{bclogo}[logo=\bcinfo, couleurBarre=orange, couleur=lightgray]{ note}}
  {\end{bclogo}}
\newenvironment{attention}
  {\begin{bclogo}[logo=\bcattention, couleurBarre=red, noborder=true, 
                couleur=red!15]{Important!}}
  {\end{bclogo}}

\begin{question}
given the function \(f(x,y,z)\)\\[0pt]
\begin{center}
\begin{tabular}{rrrrll}
x & y & z & \(f(x,y,z)\) & minterm & maxterm\\[0pt]
\hline
0 & 0 & 0 & 1 & \(x'y'z'\) & \\[0pt]
0 & 0 & 1 & 1 & \(x'y'z\) & \\[0pt]
0 & 1 & 0 & 0 &  & \(x+y'+z\)\\[0pt]
0 & 1 & 1 & 1 & \(x'yz\) & \\[0pt]
1 & 0 & 0 & 0 &  & \(x'+y+z\)\\[0pt]
1 & 0 & 1 & 0 &  & \(x'+y+z'\)\\[0pt]
1 & 1 & 0 & 1 & \(xyz'\) & \\[0pt]
1 & 1 & 1 & 0 &  & \(x'+y'+z'\)\\[0pt]
\end{tabular}
\end{center}
\begin{subquestion}
find a boolean expression for the function\\[0pt]
\begin{answer}
\begin{gather*}
  f(x,y,z) = x'y'z'+x'y'z+x'yz+xyz'\\
  f(x,y,z) = (x+y'+z)(x'+y+z)(x'+y+z')(x'+y'+z')
\end{gather*}
\end{answer}
\end{subquestion}
\begin{subquestion}
check whether \(f(x,y,z)\) is \href{boolean_algebra.org}{functionally complete}\\[0pt]
\begin{answer}
\begin{gather*}
  f(x,y,z) = x'y'z'+x'y'z+x'yz+xyz'\\
  f(x,x,x) = x'x'x' + x'x'x + x'xx + xxx' = x' + 0 + 0 + 0 = x'\\
  f(x',0,z) = x''0z' + x''0'z' + x''0z + x'0z = 0 + xz + 0 + 0 = xz
\end{gather*}
negation is derived from the second equation and the operator \(AND\) is derived from the second\\[0pt]
\(f(x,y,z)\) is functionally complete\\[0pt]
\end{answer}
\end{subquestion}
\end{question}

\begin{question}
check whether the following systems are functionally complete\\[0pt]
\begin{subquestion}
\[
f_1(x,y) = x \oplus y \qquad f_2(x,y) = x + y'
\]\\[0pt]
\begin{answer}
\(f_1(x,y)=x'y+xy'\)\\[0pt]
\(f_1(x,1)=x' \longrightarrow\) derived negation\\[0pt]
\(f_2(x, f_2(0, y)) = x+y \longrightarrow\) derived OR operator\\[0pt]
this system partially complete since i used 0/1 (using the set of operands \(\lnot\), \(\lor\))\\[0pt]
\end{answer}
\end{subquestion}
\begin{subquestion}
\[
f_1(a,b,c) = a(b+c) \qquad f_2(a) = a'
\]\\[0pt]
\begin{answer}
\(f_2(a) = a' \longrightarrow\) derived negation\\[0pt]
\(f_1(a,b,a')=a(b+a')=ab+aa'=ab+0=ab \longrightarrow\) derived AND operator\\[0pt]
since we derived both negation and the AND operator (\(\lnot, \land\)) using the system then its functionally complete\\[0pt]
\end{answer}
\end{subquestion}
\begin{subquestion}
\[
f_1(a,b) = a \oplus b \qquad f_2(a,b) = a'b
\]\\[0pt]
\begin{answer}
\(f_1(a,b)=a'b+ab'\)\\[0pt]
\(f_2(a,b)=a'b\)\\[0pt]
not a functionally complete system\\[0pt]
\end{answer}
\end{subquestion}
\end{question}

\begin{question}
write the truth table of \(f(x_1,x_2,x_3)\) which receives the value 1 if \(x_1=x_2'+x_3\)\\[0pt]
\begin{answer}
\begin{center}
\begin{tabular}{rrrrll}
\(x_1\) & \(x_2\) & \(x_3\) & \(f(x, x_2, x_3)\) & minterm & maxterm\\[0pt]
\hline
0 & 0 & 0 & 0 &  & \(x+y+z\)\\[0pt]
0 & 0 & 1 & 0 &  & \(x+y+z'\)\\[0pt]
0 & 1 & 0 & 0 &  & \(x+y'+z\)\\[0pt]
0 & 1 & 1 & 0 &  & \(x+y'+z'\)\\[0pt]
1 & 0 & 0 & 0 &  & \(x'+y+z\)\\[0pt]
1 & 0 & 1 & 1 & \(xy'z\) & \\[0pt]
1 & 1 & 0 & 0 &  & \(x'+y'+z\)\\[0pt]
1 & 1 & 1 & 0 &  & \(+y'+z'\)\\[0pt]
 &  &  &  &  & \\[0pt]
\end{tabular}
\end{center}
\end{answer}
\begin{subquestion}
find a boolean expression for the function \(f\)\\[0pt]
\begin{answer}
\[f(x,y,z) = xy'z\]\\[0pt]
\end{answer}
\end{subquestion}
\begin{subquestion}
prove that \(f\) is \href{boolean_algebra.org}{partially functionally complete}\\[0pt]
\begin{answer}
\(f(x,y,z) = xy'z\)\\[0pt]
\(f(1,y,1) = 1y'1 = y' \longrightarrow\) derived negation\\[0pt]
\(f(1,y',z) = 1yz = yz \longrightarrow\) derived the AND operator\\[0pt]
the function is partially completely functional because it can derive the already known completely functional set of operators \(\lnot, \land\)\\[0pt]
\end{answer}
\end{subquestion}
\begin{subquestion}
if \(f\) functionally complete? explain without proof\\[0pt]
\begin{answer}
\(f(x,x,x)=xxx=x\neq x'\)\\[0pt]
\(f\) isnt functionally complete\\[0pt]
\end{answer}
\end{subquestion}
\end{question}

\begin{question}
given the function \(f(x,y,z)\) which outputs the value 1 only when most of the variables are equal to 1\\[0pt]
\begin{subquestion}
write a truth table for \(f\)\\[0pt]
\begin{answer}
\begin{center}
\begin{tabular}{rrrrll}
\(x_1\) & \(x_2\) & \(x_3\) & \(f(x,x_2,x_3)\) & minterm & maxterm\\[0pt]
\hline
0 & 0 & 0 & 0 &  & \(x+y+z\)\\[0pt]
0 & 0 & 1 & 0 &  & \(x+y+z'\)\\[0pt]
0 & 1 & 0 & 0 &  & \(x+y'+z\)\\[0pt]
0 & 1 & 1 & 1 & \(x'yz\) & \\[0pt]
1 & 0 & 0 & 0 &  & \(x'+y+z\)\\[0pt]
1 & 0 & 1 & 1 & \(xy'z\) & \\[0pt]
1 & 1 & 0 & 1 & \(xyz'\) & \\[0pt]
1 & 1 & 1 & 1 & \(xyz\) & \\[0pt]
\end{tabular}
\end{center}
\end{answer}
\end{subquestion}
\begin{subquestion}
prove that \(f,\lnot\) is a partially complete system\\[0pt]
\begin{answer}
\(f(x,y,z)=x'yz+xy'z+xyz'+xyz\)\\[0pt]
\(f(x,y,z)=x'yz+xy'z+xy\)\\[0pt]
\(f(x,y,0)=0+0+xy\)\\[0pt]
\(f(x,y,0)=xy \longrightarrow\) derived the AND operator\\[0pt]
\(f\) can derive AND so we get the set \(\lnot, \land\) which is known to provide complete functionality\\[0pt]
\end{answer}
\end{subquestion}
\end{question}

\begin{question}
prove that NOR is a completely functional system\\[0pt]
\begin{answer}
\begin{center}
\begin{tabular}{rrrll}
\(x\) & \(y\) & output & minterm & maxterm\\[0pt]
\hline
0 & 0 & 1 & \(x'y'\) & \\[0pt]
0 & 1 & 0 &  & \(x+y'\)\\[0pt]
1 & 0 & 0 &  & \(x'+y\)\\[0pt]
1 & 1 & 0 &  & \(x'+y'\)\\[0pt]
\end{tabular}
\end{center}
\(f(x,y)=x'y'\)\\[0pt]
\(f(x,x)=x'x'=x' \longrightarrow\) derived negation (\(\lnot\))\\[0pt]
\(f(y,y)=y'y'=y'\)\\[0pt]
\(f(x',y')=xy \longrightarrow\) derived AND operator (\(\land\))\\[0pt]
we derived the set of operators (\(\land\), \(\lnot\)) using this function so its completely functional\\[0pt]
\end{answer}
\begin{subquestion}
build a boolean expression for \(t(x,y,z) = xy+z'\) using only NOR\\[0pt]
\begin{answer}
\[t(x,y,z)=NOR(x',y')+NOR(z,z)\]\\[0pt]
\end{answer}
\end{subquestion}
\end{question}
\end{document}